\section{El Launch File}

\subsection{\texttt{generate\_launch\_description()}}

Esta función es el punto de entrada que \texttt{ros2launch} busca en el archivo importado. Debe retornar un objeto \texttt{LaunchDescription}: una lista ordenada de acciones que el sistema de launch procesará en secuencia.

\begin{lstlisting}[style=python, caption={Estructura típica de \texttt{generate\_launch\_description()}.}]
def generate_launch_description():
    return LaunchDescription([
        DeclareLaunchArgument(
            'n_rays_used',
            default_value='0',
            description='Numero de rayos LiDAR activos'
        ),
        DeclareLaunchArgument(
            'vmax',
            default_value='1.0',
            description='Velocidad lineal maxima (m/s)'
        ),
        OpaqueFunction(function=launch_setup)
    ])
\end{lstlisting}

\subsection{\texttt{DeclareLaunchArgument}}

Esta acción registra que el launch file acepta un argumento con ese nombre, haciéndolo visible para herramientas como \texttt{ros2 launch --show-args}. Al mismo tiempo establece su valor: si el argumento llegó desde la terminal mediante \texttt{:=}, lo usa; si no llegó, aplica el \texttt{default\_value}. En ambos casos el valor permanece como string. \texttt{DeclareLaunchArgument} no evalúa ni convierte el valor bajo ninguna circunstancia. La conversión de tipo es responsabilidad del código que lo consume.

\subsection{El objeto \texttt{context}}

El \texttt{context} es el objeto central del sistema de launch. Es una estructura de datos que ROS~2 mantiene durante toda la ejecución y que contiene los valores resueltos de todos los argumentos declarados, el estado global del launch (qué procesos están activos, qué acciones han sido ejecutadas), y el mecanismo de resolución de Substitutions. No es un objeto que el usuario crea directamente; ROS~2 lo construye y lo pasa a las funciones que lo necesitan.

\subsection{\texttt{OpaqueFunction}: lógica Python arbitraria}

\texttt{OpaqueFunction} le indica al sistema de launch que, cuando llegue a esa posición de la lista, llame a la función Python especificada con el context actual. Es el mecanismo que permite ejecutar lógica arbitraria (bucles, cálculos, condicionales) dentro del launch file, algo que las acciones declarativas estándar no permiten. La función recibe el context y debe retornar una lista de acciones adicionales.

\subsection{\texttt{LaunchConfiguration} y Substitutions}

\begin{lstlisting}[style=python, caption={Resolución de argumentos dentro de \texttt{launch\_setup}.}]
def launch_setup(context, *args, **kwargs):
    n_rays = int(
        LaunchConfiguration('n_rays_used').perform(context)
    )
    vmax = float(
        LaunchConfiguration('vmax').perform(context)
    )
\end{lstlisting}

\texttt{LaunchConfiguration('n\_rays\_used')} no lee el valor inmediatamente. Crea un objeto Substitution: una promesa diferida que sabe cómo resolver su valor cuando se le proporcione un context. Esta distinción es relevante porque en el sistema de launch algunas acciones se construyen antes de que todos los argumentos estén resueltos.

\texttt{.perform(context)} resuelve esa promesa: consulta el context, recupera el string \texttt{'5'} y lo retorna. La conversión \texttt{int(...)} es responsabilidad del usuario, ya que ROS~2 nunca convierte tipos.

\subsection{Creación del descriptor de nodo}

\begin{lstlisting}[style=python, caption={Creación del descriptor de nodo con parámetros.}]
Node(
    package='control_nodes',
    executable='navigate_individual_pva',
    namespace='robot0',
    parameters=[{
        'n_rays_used': n_rays,  # int 5
        'vmax':        vmax,    # float 0.8
    }],
)
\end{lstlisting}

El objeto \texttt{Node} no lanza el proceso todavía. Es un descriptor que el sistema de launch procesará cuando llegue su turno. Al procesarlo, serializa el diccionario \texttt{parameters} a YAML y lo pasa al proceso como argumento de línea de comandos. Lo que el sistema operativo termina ejecutando es equivalente a:

\begin{lstlisting}[style=bash, caption={Comando real que el OS ejecuta al procesar el descriptor \texttt{Node}.}]
navigate_individual_pva \
    --ros-args \
    -p n_rays_used:=5 \
    -p vmax:=0.8 \
    -r __ns:=/robot0
\end{lstlisting}

El diccionario Python se convirtió en flags de CLI. Esta es la segunda frontera de serialización del flujo: Python dict $\rightarrow$ YAML $\rightarrow$ argumentos de proceso.
